% =====================================================================
% Reemplaza la subsección «El Concepto de Campo Gravitacional ($\boldsymbol{g}$)»,
% que sigue a la nueva subsección «La ley de gravitación universal en
% términos del campo» (archivo gravitacion_campo.tex).
% La etiqueta eq:campo_g ahora está en gravitacion_campo.tex; aquí ya no
% se define, para no duplicarla. Se conserva fig:campo_gravitacional_newton.
% Requiere: tikz con \usetikzlibrary{arrows.meta}.
% =====================================================================

\subsection{Líneas de campo y superficies equipotenciales}
\label{subsec:lineas_equipotenciales}

Newton se sentía incómodo con la idea de una «acción a distancia», en la
que un cuerpo tira de otro a través del espacio vacío. El concepto de campo
evita esa dificultad: la masa $M$ modifica el espacio a su alrededor, y una
masa $m$ responde al campo que encuentra \emph{en el lugar donde está},
mediante $\boldsymbol{F} = m\,\boldsymbol{g}$. El campo es una propiedad del
espacio creada por $M$, haya o no otra masa que lo sienta.

La figura~\ref{fig:campo_gravitacional_newton} muestra el campo de la
ecuación~\eqref{eq:campo_g_vectorial} alrededor de la Tierra. Las
\textbf{líneas de campo} son radiales y apuntan hacia el centro; donde están
más juntas, el campo es más intenso. Las \textbf{superficies equipotenciales}
son esferas centradas en la Tierra: sobre cada una el campo tiene la misma
magnitud, y en todos sus puntos es perpendicular a ella. Su significado en
términos de energía se verá al estudiar la energía potencial gravitacional.

\begin{figure}[htbp]
  \centering
  \begin{tikzpicture}[
      scale=1.2, font=\large,
      tierra/.style={shade, ball color=blue!60!cyan, text=white},
      equipotencial/.style={dashed, orange, thick},
      lineacampo/.style={-{Stealth[length=3mm, width=2mm]}, ultra thick, red!80!black},
    ]
    % Tierra en el centro
    \node[tierra, circle, minimum size=3cm] at (0,0) {Tierra};
    % Superficies equipotenciales (cortes de esferas)
    \foreach \r in {2.5, 3.5, 4.5} {
      \draw[equipotencial] (0,0) circle (\r);
    }
    % Líneas del campo gravitacional
    \foreach \ang in {0, 45, ..., 315} {
      \draw[lineacampo] (\ang:5) -- (\ang:1.5);
    }
    % Rótulos
    \node[orange, font=\small, anchor=south west] at (0.1, 3.5) {Equipotencial};
    \node[red!80!black, font=\small, rotate=45, anchor=south]
      at (45:3.5) {Campo gravitacional $\boldsymbol{g}$};
  \end{tikzpicture}
  \caption{Corte en dos dimensiones del campo gravitacional $\boldsymbol{g}$
    de la Tierra. Las líneas de campo (flechas rojas) son radiales y apuntan
    hacia el centro; las equipotenciales (círculos naranjas a trazos) son
    perpendiculares a ellas.}
  \label{fig:campo_gravitacional_newton}
\end{figure}

\paragraph{El campo cerca de la superficie.}
Para alturas $h$ pequeñas comparadas con el radio terrestre, $h \ll R_T$, la
distancia al centro casi no cambia, $r = R_T + h \approx R_T$, y en una
región pequeña las líneas de campo son prácticamente paralelas. El campo es
entonces casi uniforme y, con el eje $y$ vertical hacia arriba,
\begin{equation}
  \boldsymbol{g} \approx -\,g_0\,\hat{\boldsymbol{y}} \,,
  \qquad
  g_0 = \frac{G M_T}{R_T^{2}} \approx \SI{9.8}{\meter\per\second\squared} \,.
  \label{eq:campo_uniforme}
\end{equation}
Este es el campo que usamos en los problemas de caída libre y de movimiento
parabólico. Por otra parte, las leyes de Kepler, que describen el movimiento
de los planetas, se deducen de que la gravitación es una interacción central
que decae como el inverso del cuadrado de la distancia.
